% !TeX program = pdflatex
% Standalone document: compile twice with pdflatex, or use latexmk -pdf.
% No external figures, bibliography database, shell escape, or code execution.
% Historical proposal. The user subsequently authorized terminology migration.
% See docs/migration/UVLM_ANALYSIS_GUIDE.tex for implemented interfaces and limits.
\documentclass[11pt,a4paper]{article}
\usepackage[T1]{fontenc}
\usepackage[utf8]{inputenc}
\usepackage{lmodern}
\usepackage[margin=24mm,headheight=14pt]{geometry}
\usepackage{microtype}
\usepackage{amsmath,amssymb}
\usepackage{array,booktabs,longtable}
\usepackage{enumitem}
\usepackage{xcolor}
\usepackage{listings}
\usepackage{xurl}
\usepackage{fancyhdr}
\usepackage{hyperref}
\definecolor{ink}{HTML}{16324F}
\definecolor{accent}{HTML}{19647E}
\definecolor{papergray}{HTML}{F4F6F8}
\hypersetup{colorlinks=true,linkcolor=ink,urlcolor=accent,
  pdftitle={UVLM to AeroBeams: Terminology and Architecture Migration Proposal},
  pdfauthor={Technical planning document},bookmarksnumbered=true}
\urlstyle{tt}
\newcommand{\code}{\nolinkurl}
\newcommand{\proposed}{\textbf{Proposed interface.}\ }
\newcommand{\current}{\textbf{Current implementation.}\ }
\newcolumntype{L}[1]{>{\raggedright\arraybackslash}p{#1}}
\setlength{\parindent}{0pt}
\setlength{\parskip}{0.5em}
\setlength{\emergencystretch}{2em}
\renewcommand{\arraystretch}{1.2}
\setlist{itemsep=0.2em,topsep=0.3em,leftmargin=*}
\setcounter{tocdepth}{2}
\pagestyle{fancy}
\fancyhf{}
\fancyhead[L]{\small UVLM / AeroBeams migration}
\fancyhead[R]{\small Proposal for review}
\fancyfoot[C]{\thepage}
\lstdefinestyle{juliaplan}{
  basicstyle=\ttfamily\footnotesize,
  backgroundcolor=\color{papergray},
  frame=single,rulecolor=\color{black!15},
  breaklines=true,breakatwhitespace=false,
  columns=fullflexible,keepspaces=true,
  showstringspaces=false,tabsize=4,
  xleftmargin=4pt,xrightmargin=4pt,
  aboveskip=0.8em,belowskip=0.8em,
  escapeinside={(*@}{@*)}
}
\lstset{style=juliaplan}

\begin{document}
\begin{titlepage}
\vspace*{18mm}
{\color{accent}\Large Technical migration proposal\par}
\vspace{8mm}
{\color{ink}\Huge\bfseries From the current UVLM code\\[3mm]
to AeroBeams terminology\\[3mm]
and structure\par}
\vspace{10mm}
{\Large Meaning of the names, ownership of the data,\\
solver interfaces, and implementation sequence\par}
\vspace{15mm}
\hrule
\vspace{5mm}
\textbf{Historical proposal: terminology migration subsequently authorized.}

The delivered interfaces, retained convergence studies and remaining roadmap
items are documented in \code{docs/migration/UVLM_ANALYSIS_GUIDE.tex}.
Approval language elsewhere below records the original proposal checkpoint.

This document proposes changes. It does not authorize or implement the migration.
The first delivery would reorganize the existing UVLM interface while preserving
its numerical calculations. Actual coupling to AeroBeams would follow through
separate, validated implementation stages.
\vspace{5mm}
\hrule
\vfill
\textbf{Reference checkout}\par
AeroBeams 0.8.1; WingPropellerUVLM 0.1.0.\par
Planning baseline: \code{d0fe40bf7f58f174c49cac62c62e424e32b41586},
including the subsequent moving-block damping changes.\par
Prepared 14 September 2026.
\end{titlepage}

\tableofcontents
\clearpage

\section{Purpose and proposed decision}
\label{sec:decision}

The objective is to make the wing--propeller UVLM code recognizable to an
AeroBeams user: build a model, choose an analysis problem and solver, call
\code{solve!}, and inspect clearly named results. Achieving this requires more
than changing spelling. Some existing names describe several responsibilities
at once, and some similar names in the two libraries describe different physics.

The recommendation is to preserve \code{WingPropellerUVLM} as a standalone
aerodynamic package, adopt AeroBeams-style constructors and problem objects,
and introduce a separate proposed \code{AeroBeamsUVLM} bridge. The bridge would
transfer motion and dimensional loads while coordinating accepted time steps.
AeroBeams would remain responsible for the structural equations and structural
integration in the coupled analysis.

\subsection{Three different kinds of change}

\begin{longtable}{@{}L{0.18\linewidth}L{0.36\linewidth}L{0.38\linewidth}@{}}
\toprule
Change & Example & Required evidence\\
\midrule\endhead
Terminology & \code{maximum_iterations} becomes
\code{maximumIterations}. & Both inputs resolve to the same value, with the
same residual definition and numerical result.\\
Organization & The current \code{System} is presented through separate model,
state, and workspace objects. & No duplicated ownership, unintended array
aliasing, changed wake chronology, or numerical drift.\\
Formulation & The Chang linear structural integrator is replaced by an
AeroBeams structural problem. & Structural and coupled physical benchmarks;
matching names alone cannot establish equivalence.\\
\bottomrule
\end{longtable}

This distinction determines the order of implementation. First preserve the
existing calculation behind a clearer interface. Then introduce the geometric
and structural bridge. Numerical extensions such as adaptive stepping or
linearized UVLM eigenanalysis follow after the basic coupling is established.

\subsection{What approval would cover}

The recommended initial approval scope is phases P0--P5: baseline capture,
environment isolation, public terminology, ownership separation, studies,
results, and provenance. Phases P6--P9 describe the subsequent AeroBeams bridge
and rotor validation. They can be approved together with the whole roadmap or
as a later implementation milestone. Advanced analyses in P11 require their
own numerical design and validation.

\textbf{Until approval is given, all names and examples labeled proposed in
this document remain design specifications.} Existing simulation files,
settings, and aerodynamic selections are not changed by preparing this plan.

\subsection{How to read this document}

Read Sections~\ref{sec:concepts}--\ref{sec:dictionary} to understand the names,
Sections~\ref{sec:semantics}--\ref{sec:lifecycle} for the main technical
differences, and Sections~\ref{sec:phases}--\ref{sec:approval} for the work and
acceptance decision. The code listings illustrate an intended API. They are
not runnable examples of interfaces that already exist.

\section{What AeroBeams means by model, problem, and solver}
\label{sec:concepts}

\subsection{Model: the physical assembly}

An AeroBeams \code{Model} describes an assembly of beams, boundary conditions,
connections, inertial properties, prescribed reference motion, and related
physical data. A \code{Beam} carries its geometry, sectional properties, and
element discretization. The public constructors include \code{create_Beam}
and \code{create_Model}.

For UVLM, the proposed \code{UVLMModel} would contain aerodynamic topology,
reference surface geometry, rotor definitions, surface IDs, and reference
quantities for coefficients. Creating it would validate the geometry and build
derived topology. It would not simulate six seconds of motion, create plots,
or write a convergence result.

\textbf{Practical example.} Increasing airspeed from 65 to 85~m/s generally
creates a new operating point for the same reference geometry. Increasing the
number of blade panels changes the aerodynamic model discretization. These
operations should not both be implemented as arbitrary mutations of one
large run object.

\subsection{Problem: the analysis to perform}

An AeroBeams \code{DynamicProblem} combines a model, a system solver, time
settings, the current numerical state, and tracked output. Its existing
constructor uses names such as \code{initialTime}, \code{finalTime},
\code{timeVector}, \code{trackingTimeSteps}, and \code{trackingFrequency}.
The time-step keyword is the mathematical symbol $\Delta t$ in the Julia API.

The corresponding standalone aerodynamic object would be a
\code{UVLMDynamicProblem}. It would answer: which model, at which operating
point, over which time grid, with which initial wake and output policy?
A separate \code{UVLMSteadyProblem} would describe a supported steady
aerodynamic calculation.

The proposed coupled wrapper, \code{CoupledDynamicProblem}, would contain a
real AeroBeams \code{DynamicProblem} and references to the UVLM model, state,
workspace, and coupling configuration. It would not disguise the old Chang
state vector as an AeroBeams structural problem.

\subsection{Solver: the algorithm and convergence controls}

AeroBeams uses \code{create_NewtonRaphson} to configure the structural nonlinear
solution. The UVLM migration needs two additional concepts:

\begin{itemize}
\item \code{UVLMSolver}: aerodynamic algorithm choices, including force model,
core policy, interactions, wake treatment, and applicable matrix reuse.
\item \code{PartitionedCoupling}: outer iteration controls, including state
and load tolerances, maximum iterations, and relaxation.
\end{itemize}

The structural Newton solver and the outer partitioned solver have different
residuals. Setting both fields to the same numerical tolerance does not make
their convergence criteria equivalent.

\subsection{State, workspace, and results}

\begin{longtable}{@{}L{0.24\linewidth}L{0.34\linewidth}L{0.34\linewidth}@{}}
\toprule
Proposed object & Question it answers & Typical contents\\
\midrule\endhead
\code{UVLMState} & What aerodynamic history has been accepted? & Bound/wake
circulation, wake positions, previous geometry, active wake rows, rotor phase.\\
\code{UVLMWorkspace} & What storage is needed to compute the next trial? & AIC,
right-hand side, factorization cache, temporary velocities, force buffers.\\
\code{UVLMResults} & What should the user inspect or save? & Accepted time
coordinates, response histories, loads, status, optional geometry snapshots.\\
\bottomrule
\end{longtable}

A workspace is disposable scratch storage. A state contains history that
affects the next solution. A result is a saved record. A saved result must not
change when a later trial overwrites a workspace array.

These UVLM object names are proposed organizational counterparts. AeroBeams
does not already provide a one-to-one set of these UVLM types. Compatibility
means matching its conventions and lifecycle where meaningful, not claiming
that both libraries have identical internals.

\section{Naming rules and a detailed terminology dictionary}
\label{sec:dictionary}

\subsection{Rules to apply consistently}

\begin{enumerate}
\item Use PascalCase for public types: \code{UVLMModel}, \code{UVLMState}.
\item Use keyword constructors named \code{create_Type}, following
\code{create_Model} and \code{create_DynamicProblem}.
\item Use snake_case operation names, with \code{!} for mutation:
\code{update_geometry!}, \code{commit_time_step!}.
\item Use camelCase for descriptive public fields and keywords:
\code{maximumIterations}, \code{nChordwisePanels}.
\item Keep established mathematical symbols when their meaning is explicit.
Do not replace every internal short variable with a long descriptive name.
\item Document units, basis, ownership, and whether a value is prescribed,
derived, trial, accepted, or saved.
\end{enumerate}

Keep the UVLM prefix on exported model/problem types to avoid ambiguity when
both packages are loaded. Internal backend routines can retain their original
names during the first migration. Using concrete or parameterized storage is
preferable in hot loops; broad \code{Real} or \code{Function} fields do not
need to be copied merely because some AeroBeams types use them.

\subsection{Objects and operation names}

\begingroup\small
\begin{longtable}{@{}L{0.29\linewidth}L{0.29\linewidth}L{0.34\linewidth}@{}}
\toprule
Current UVLM concept & Proposed terminology & Meaning and migration action\\
\midrule\endhead
\code{System} & \code{UVLMModel}, \code{UVLMState}, \code{UVLMWorkspace} &
Separate physical definition, numerical history, and scratch arrays. Keep a
backend facade during transition.\\
\code{initialize_bohnisch_uvlm_system} & \code{create_UVLMModel} and
problem initialization & Move Bohnisch-specific defaults into a named case
builder. Generic construction must not encode a particular experiment.\\
\code{ChangModel}, \code{build_chang_model} & Benchmark
\code{create_ChangModel} & Preserve the current linear benchmark; later add a
separate AeroBeams-based builder.\\
\code{run_chang} & Construction, \code{solve!}, explicit result writing &
Retain a compatibility wrapper so current workflows keep working.\\
\code{steady_analysis!} & \code{solve!} on \code{UVLMSteadyProblem} &
Wrap the existing backend before changing numerical behavior.\\
\code{unsteady_analysis!} & \code{solve!} on \code{UVLMDynamicProblem} &
Distinguish the full analysis from a single physical-step trial.\\
\code{PartitionedCouplingOptions} & \code{create_PartitionedCoupling} &
Provide a validated public constructor; preserve outer residual definitions.\\
\code{snapshot_uvlm}, \code{restore_uvlm!} & Explicit step snapshot and
\code{rollback_time_step!} & Restore all history needed for a deterministic
trial, not just circulation.\\
\code{advance_uvlm_trial!} & \code{evaluate_trial!} & A coupling trial must
not advance the accepted wake by default.\\
\code{advance_wake!}, \code{commit_wake_rows!} &
\code{commit_time_step!} orchestration & Advance wake/history and active-row
counts exactly once after acceptance.\\
\bottomrule
\end{longtable}
\endgroup

\subsection{Geometry, operating point, and solver settings}

\begingroup\small
\begin{longtable}{@{}L{0.29\linewidth}L{0.29\linewidth}L{0.34\linewidth}@{}}
\toprule
Current field & Proposed field & What must be preserved\\
\midrule\endhead
\code{speed_mps}, \code{freestream_speed_mps} & \code{airspeed} &
m/s; also document the full freestream vector and basis.\\
\code{rpm}, \code{rotation_rpm} & Input \code{rotationRPM}; resolved
\code{angularSpeed} & Convert RPM to rad/s once.\\
Proportional speed override & \code{rotorSpeedPolicy} with
\code{:constantAdvanceRatio} & Scale each prescribed rotor speed from its
reference point and keep that reference in metadata.\\
\code{angle_of_attack_deg}, \code{sideslip_deg} & \code{angleOfAttack},
\code{sideslip} & Resolved values are radians; legacy degrees require conversion.\\
\code{azimuth_step_deg}, \code{azimuth_deg} & \code{azimuthStep} &
Radians; derive physical time step from the intended rotor speed policy.\\
\code{spanwise_panels}, \code{ns} & \code{nSpanwisePanels} &
Aerodynamic discretization, independent of structural \code{nElements}.\\
\code{radial_panels} & \code{nRadialPanels} & Number of blade radial panels.\\
\code{chordwise_panels}, \code{nc} & \code{nChordwisePanels} &
Keep the backend array ordering unchanged initially.\\
Maximum wake rows, \code{nw} & \code{maximumWakeRows} & Allocated capacity.\\
\code{nwake}, \code{iwake} & \code{activeWakeRows} & Accepted/trial active
length, bounded by capacity.\\
\code{wake_revolutions} & \code{wakeRevolutions} & Retention policy;
conversion to rows depends on the time/azimuth grid.\\
\code{core_radius_m} & \code{coreRadius} & Fixed physical radius in metres.\\
\code{segment_core_factor}, \code{chord_core_factor} &
\code{segmentCoreFactor}, \code{chordCoreFactor} & Keep the factor-based
core law distinct from a fixed-radius law.\\
\code{interaction_on}, \code{surface_id} & Resolved
\code{interactionGroups} policy & Disabling cross-group interaction must
retain interaction between blades of one rotor.\\
\code{attachment_eta} & \code{attachments} & Store beam identity and
physical/normalized attachment coordinate; interpolate off-node attachment.\\
\code{elastic_axis_fraction} & \code{normSparPos}, if equivalent &
Verify chord origin, direction and offset before unifying the names.\\
\code{maximum_iterations} & \code{maximumIterations} & Same algorithmic limit.\\
\code{state_tolerance}, \code{load_tolerance} & \code{stateTolerance},
\code{loadTolerance} & Preserve scaling, denominator and tested quantity.\\
\code{equilibrium_tolerance} & \code{equilibriumTolerance} &
Do not assume that the Chang and AeroBeams equilibrium residuals are identical.\\
\bottomrule
\end{longtable}
\endgroup

\subsection{Outputs, damping, and initialization}

\begingroup\small
\begin{longtable}{@{}L{0.29\linewidth}L{0.29\linewidth}L{0.34\linewidth}@{}}
\toprule
Current concept & Proposed concept & Important distinction\\
\midrule\endhead
\code{surface_history} & Optional \code{surfaceGeometryOverTime} & Heavy
geometry history is independent of dense lightweight response samples.\\
Response \code{time} & \code{responseTimeVector} & The time coordinates
must match the exact response samples used by the estimator.\\
Tracked states & \code{savedTimeVector} and named histories & Match the
AeroBeams tracking vocabulary, documenting each sampling stride.\\
\code{moving_block_lambda_per_s} & \code{growthRate} &
$\lambda$ in s$^{-1}$; negative means decay.\\
\code{frequency_hz} & \code{frequencyHz} & Cycles per second.\\
\code{omega_radps} & \code{angularFrequency} & Radians per second.\\
\code{damping_ratio_percent} & \code{dampingRatio}, explicit percent export &
Internal dimensionless fraction and plotted percentage are different.\\
Startup \code{trim_revolutions} & \code{baselineRevolutions} &
Mean aerodynamic load preparation about the current baseline; not a solved
deformed structural trim.\\
\code{validation_passed} & Separate \code{completed},
\code{solverConverged}, \code{acceptedForStudy} & Distinguish finishing the
calculation, satisfying equations, and satisfying study criteria.\\
\bottomrule
\end{longtable}
\endgroup

Julia field names and serialized keys have different audiences. Changing a
public Julia field to camelCase need not break existing snake_case CSV/TOML
columns. Preserve readers and version any necessary output changes.

\section{Main semantic and numerical differences}
\label{sec:semantics}

\subsection{The structural state vectors are not interchangeable}

The Chang benchmark advances a linear structural coordinate vector through
mass, damping/gyroscopic, and stiffness matrices. Conceptually its equations
have the familiar form
\begin{equation}
 M\ddot q+C\dot q+Kq=Q_{\mathrm{aero}}+Q_{\mathrm{excitation}}.
\end{equation}
The exact contents and constraints of $q$ belong to that benchmark.

AeroBeams uses a mixed geometrically exact beam formulation. Its elemental
states include
\begin{equation}
 \{u,p,F,M,V,\Omega,\chi\},
\end{equation}
where $u$ and $p$ are deformation variables, $F$ and $M$ are sectional
resultants, $V$ and $\Omega$ are sectional velocities, and $\chi$ denotes
the existing local aerodynamic states when present. The global unknown
vector also reflects constraints and formulation-specific degrees of freedom.

Consequently, a Chang expression that extracts six coordinates per node
cannot be reused by changing the variable name to \code{problem.x}. The bridge
must obtain element/node states using AeroBeams IDs and reconstruction logic.
Structural-only mode frequencies, static compliance, mass, and attachment
motion must be compared before attempting aeroelastic parity.

\subsection{Euler angles and Wiener--Milenkovic parameters}

The active Chang adapter constructs rotations using its Euler-angle
composition. AeroBeams deformation rotations use Wiener--Milenkovic
parameters. A vector of Euler angles is therefore not the same quantity as
AeroBeams' $p$, even if both have three components.

For the bridge, reconstruct the deformation rotation through the existing
\code{rotation_tensor_WM} utilities and compose it with the initial section
orientation consistently. The initial beam \code{rotationParametrization}
is a separate issue from the representation of its deformation state.

The rate $\dot p$ is not generally the physical angular velocity $\omega$.
They are related by a parameterization-dependent tangent map. The same issue
appears when converting a generalized rotational force into a physical moment.
This is why renaming Euler coordinates to $p$ would introduce a physical error.

\subsection{Panel count and beam-element count}

The current Chang setup ties wing spanwise aerodynamic panels to structural
elements. The proposed interface exposes these independently:
\begin{equation}
 n_{\mathrm{spanwise\ panels}}\quad\hbox{and}\quad
 n_{\mathrm{structural\ elements}}.
\end{equation}
A beam might use 20 elements while the aerodynamic wing uses 30 spanwise
panels and 10 chordwise panels. The bridge interpolates the beam motion to
aerodynamic vertices and transfers loads back consistently. These numbers
are illustrative, not recommended converged values.

An attachment at 83\% of the modeled span must remain at that physical
position after either mesh is refined. Rounding it to the nearest beam node
would change the modeled installation during a convergence study.

\subsection{Local aerodynamics and global UVLM}

Existing AeroBeams aerodynamic models are evaluated through element-local
aerodynamic properties and states. UVLM solves a circulation system whose
surfaces and wakes interact globally. Its schematic aerodynamic solve is
\begin{equation}
 A(\mathcal{G},\mathcal{W})\,\Gamma=b(\mathcal{G},\mathcal{W},v),
\end{equation}
where $\mathcal{G}$ denotes the current geometry and $\mathcal{W}$ the
retained wake/history.

A separate UVLM solve inside every element's \code{aerodynamic_loads!}
would repeat the global calculation and could advance the same wake multiple
times. The bridge must evaluate the global aerodynamic subsystem once for
each coupling trial and distribute its loads to the participating structure.
Surfaces must not simultaneously receive full sectional and full UVLM loads
unless a deliberately defined combination model is being used.

\subsection{Dimensional loads and aerodynamic coefficients}

Structural equilibrium requires forces in N and moments in N\,m. UVLM
coefficient-output functions may instead normalize loads by reference density,
velocity, area, or length. Coefficients are useful for plots and convergence
metrics, but they are not structural load vectors.

The interface must document every stored circulation and force quantity.
In particular, the backend's freestream derivative storage must not be
interpreted as a temporal increment merely because its name begins with ``d''.
Audit the source fields denoted $\mathrm{d}\Gamma$ and
$\mathrm{d}\Gamma\mathrm{dt}$ when implementing the contract; the first
contains freestream derivatives and the second a temporal rate.

AeroBeams' \code{forceScaling} is numerical residual scaling. The load buffer
should remain dimensional, and scaling should be applied once at residual
assembly. The current \code{UnitsSystem} labels affect plotting; they do not
convert numerical input values.

\subsection{Growth rate, angular frequency, and damping ratio}
\label{sec:dampingmeaning}

For a decaying or growing single-mode response,
\begin{equation}
 x(t)\simeq A_0 e^{\lambda t}\sin(\omega_d t+\varphi),
 \qquad \omega_d=2\pi f,
\end{equation}
$\lambda$ is the exponential growth rate, $f$ is frequency in Hz, and
$\omega_d$ is angular frequency in rad/s. The associated modal damping ratio is
\begin{equation}
 \zeta=-\frac{\lambda}{\sqrt{\lambda^2+\omega_d^2}}.
\end{equation}
Negative $\lambda$ means decay and positive $\zeta$; positive $\lambda$
means growth and negative $\zeta$.

The local AeroBeams eigenproblem stores \code{frequencies} as angular
frequencies and \code{dampings} as real eigenvalue parts. For a matched
oscillatory mode, the latter has the growth-rate sign convention above. It
is not already a percentage damping ratio.

For example, $f=4$~Hz and $\lambda=-0.15$~s$^{-1}$ correspond to
$\omega_d\approx25.13$~rad/s and $\zeta\approx0.00597$, or about
$0.597\%$. Comparing ``4'' directly to an AeroBeams frequency of ``25.13''
would be a unit error, not a disagreement in the physical mode.

\subsection{Baseline subtraction, steady equilibrium, and trim}

The current startup procedure collects mean aerodynamic loads about the
chosen configuration and subtracts that baseline during the response. It
does not automatically solve a deformed equilibrium configuration.

AeroBeams steady analysis solves a steady equilibrium; trim also introduces
appropriate unknown operating/control variables. Therefore a field such as
\code{trim_revolutions} should become \code{baselineRevolutions} in the
present workflow. A future \code{TrimProblem} must implement actual trim
semantics instead of inheriting that name from the startup procedure.

With rotating propellers, the baseline may be periodic. A phase-averaged
equilibrium, an azimuth-frozen calculation, and a periodic orbit are different
approximations. The selected interpretation must be stated before adding
eigenanalysis or interpreting a moving-block fit as free modal damping.

\subsection{Structural integration and physical damping}

The current Chang calculation uses its generalized-alpha structural scheme.
AeroBeams dynamic integration computes rates within its own problem lifecycle.
Once AeroBeams supplies the structure, the bridge must not also integrate the
same structural state with the Chang generalized-alpha corrector.

The local AeroBeams version does not provide general structural damping as
a direct counterpart to the current Chang damping-matrix options. A nonzero
pylon/mount damping specification requires a documented physical component
or a separately supported legacy model. It cannot silently become zero, and
numerical time-integration dissipation is not a physical substitute.

\section{Target architecture and data ownership}
\label{sec:architecture}

\subsection{Three packages with one dependency direction}

\begin{center}
\setlength{\fboxsep}{8pt}
\fbox{\begin{minipage}{0.84\linewidth}\centering
\textbf{AeroBeamsUVLM: proposed bridge}\\[3pt]
Mapping, coupled problem, trial/commit coordination\\[6pt]
$\swarrow\hspace{45mm}\searrow$\\[5pt]
\textbf{AeroBeams}\hfill\textbf{WingPropellerUVLM}\\[3pt]
Structure and constraints\hfill Aerodynamics and wake
\end{minipage}}
\end{center}

The bridge depends on both discipline packages. Neither discipline package
depends on the bridge. This avoids a circular dependency and allows rigid
aerodynamic studies without loading AeroBeams.

Small generic structural facilities are still needed in AeroBeams: external
dimensional resultants and safe single-step preparation/acceptance/rollback.
These facilities should not contain UVLM equations. This plan selects a
separate bridge package rather than leaving the choice between a bridge and
an extension unresolved.

\subsection{Proposed file organization}

\begin{lstlisting}
lib/WingPropellerUVLM/
  src/WingPropellerUVLM.jl
  src/OperatingPoint.jl
  src/UVLMSurface.jl
  src/UVLMModel.jl
  src/UVLMSolver.jl
  src/UVLMState.jl
  src/UVLMWorkspace.jl
  src/UVLMProblem.jl
  src/UVLMResults.jl
  src/Core.jl
  src/Kinematics.jl
  src/Loads.jl
  src/Compatibility.jl
  src/backend/             # retained numerical kernels
  src/aeroelastic/         # retained linear benchmark
  studies/Project.toml
  studies/src/Convergence.jl
  studies/src/MovingBlockDamping.jl
  studies/src/Provenance.jl
  examples/chang_linear_aeroelastic/
  test/

lib/AeroBeamsUVLM/
  Project.toml
  src/AeroBeamsUVLM.jl
  src/CouplingFrames.jl
  src/BeamUVLMCoupling.jl
  src/StructuralKinematics.jl
  src/LoadTransfer.jl
  src/PartitionedCoupling.jl
  src/CoupledDynamicProblem.jl
  src/SpinningRotor.jl
  examples/
  test/

# Generic additions/changes in the root AeroBeams package:
src/ExternalResultants.jl
src/Model.jl
src/Problem.jl
src/SystemSolver.jl
\end{lstlisting}

Create these files as functionality moves, not as an empty framework.
Case tables and experiment-specific defaults stay in example builders.
General convergence machinery and moving-block postprocessing move to a
declared studies environment. Plotting remains optional.

\subsection{A gradual split of the current System}

Initially wrap the current backend \code{System}; do not rewrite every kernel.
Inventory each field as model data, accepted history, trial data, scratch,
or output. Then introduce explicit ownership and migrate one group at a time.

The facade and new state must refer to one authoritative circulation/wake
storage. Keeping two arrays and occasionally copying between them creates a
second state synchronization problem. Add tests that verify array identities,
independent runs, and saved-result isolation before changing storage layout.

\section{Before-and-after workflow examples}
\label{sec:workflow}

\subsection{Current editable case and run entry}

The present workflow exposes case-specific settings and calls the Chang
driver, which builds the model/workspace, executes, and writes output.
A representative settings fragment is:

\begin{lstlisting}
simulation = (
    freestream_speed_mps = 85.0,
    azimuth_step_deg = 2.5,
    end_time_s = 6.0,
)
# Existing run_chang(config) handles several run responsibilities.
\end{lstlisting}

\subsection{Proposed standalone aerodynamic workflow}

\proposed The listing uses abstract geometry inputs already supplied by a
case builder. These constructors are targets, not existing calls.

\begin{lstlisting}
import WingPropellerUVLM as UVLM

model = UVLM.create_UVLMModel(; surfaces, rotors, reference)
solver = UVLM.create_UVLMSolver(; coreRadius, interactionGroups)
operatingPoint = UVLM.create_OperatingPoint(;
    airspeed = 85.0,
    referenceAirspeed = 65.0,
    referenceRotationRPM = referenceRPM,
    rotorSpeedPolicy = :constantAdvanceRatio,
)
problem = UVLM.create_UVLMDynamicProblem(;
    model, aeroSolver = solver, operatingPoint,
    (*@$\Delta t$@*) = timeStep, finalTime = 6.0,
    trackingTimeSteps = true,
    trackingFrequency = 1,
)
UVLM.solve!(problem)
results = problem.results
\end{lstlisting}

The constructor resolves and validates inputs. \code{solve!} mutates the
problem. Output writing is a separate operation. A retained \code{run_chang}
wrapper can still perform all three for backward compatibility.

\subsection{Proposed AeroBeams-coupled workflow}

\proposed Structural construction continues to use native AeroBeams calls;
the bridge owns the additional coupling objects.

\begin{lstlisting}
import AeroBeams
import AeroBeamsUVLM as Bridge

structuralModel = AeroBeams.create_Model(; beams, BCs)
structuralProblem = AeroBeams.create_DynamicProblem(;
    model = structuralModel,
    systemSolver = AeroBeams.create_NewtonRaphson(),
    (*@$\Delta t$@*) = timeStep, finalTime = 6.0,
)
mapping = Bridge.create_BeamUVLMCoupling(;
    structuralModel, aerodynamicModel = model,
    attachments, frames,
)
coupled = Bridge.create_CoupledDynamicProblem(;
    structuralProblem,
    aerodynamicModel = model,
    aerodynamicSolver = solver,
    operatingPoint,
    coupling = mapping,
    couplingSolver = Bridge.create_PartitionedCoupling(;
        maximumIterations = 10,
        stateTolerance = 1e-5,
        loadTolerance = 1e-2,
    ),
)
AeroBeams.solve!(coupled)
\end{lstlisting}

The bridge would explicitly extend \code{AeroBeams.solve!} for its own wrapper
type. Merely creating a subtype of AeroBeams \code{Problem} does not activate
a solver: the existing generic solve dispatch contains exact-type branches.
Qualified module calls also avoid two ambiguous unqualified \code{solve!}
exports in user scripts.

\subsection{The requested speed override remains explicit}

The user-selected constant-advance-ratio policy is
\begin{equation}
 N_{\mathrm{new}}=N_{\mathrm{ref}}\frac{U_{\mathrm{new}}}{U_{\mathrm{ref}}},
 \qquad \Omega_{\mathrm{new}}=\frac{2\pi N_{\mathrm{new}}}{60}.
\end{equation}
Here $N$ is RPM and $\Omega$ is rad/s. For 65 to 85~m/s, RPM scales by
$85/65$. At fixed azimuth increment,
\begin{equation}
 \Delta t=\frac{\Delta\psi}{|\Omega|},
\end{equation}
with $\Delta\psi$ in radians. Thus changing speed also changes the time
step and wake-age interpretation if azimuth resolution is held fixed.
Zero spin needs a separately specified physical time step.

The original aerodynamic selection remains the reference record. The
aeroelastic operating point is recorded separately; it must not inherit an
unchanged verification seal that claims verification at the new speed.

\section{Motion transfer and conservative load transfer}
\label{sec:transfer}

\subsection{An explicit frame convention}

Define $R_{UA}$ by
\begin{equation}
 v_U=R_{UA}v_A,\qquad R_{UA}^{T}R_{UA}=I,\qquad\det R_{UA}=+1.
\end{equation}
Document the frame origins as well as their orientation. Position mapping
includes reference beam geometry, displacement, section rotation, chord
offset, attachment offset, and rotor phase. If the reference frame moves,
include origin velocity and angular transport consistently in velocities
and wake convection.

Neither flipping a plotted axis nor relabeling the force components is a
complete coordinate transformation. A rigid rotation test should transform
both geometry and loads and preserve the physical result.

\subsection{Work conservation as the interface test}

For an aerodynamic vertex-position map $x_a=g(q_s,t)$, define the linearized
motion transfer $H$ by
\begin{equation}
 \delta x_a=H\,\delta q_s.
\end{equation}
Its work-conjugate load transfer is
\begin{equation}
 Q_s=H^T f_a,\qquad f_a^T\delta x_a=Q_s^T\delta q_s.
\end{equation}
The position associated with each force must match the position used in the
UVLM force formulation. For a moment moved between reference points,
\begin{equation}
 m_{\mathrm{new}}=m_{\mathrm{old}}+
 (r_{\mathrm{old}}-r_{\mathrm{new}})\times f.
\end{equation}

A force conjugate to a rotation parameter is not automatically the physical
moment used by the structural residual. The rotational tangent conversion
must be derived and tested. Total force and total moment conservation are
necessary checks, but the virtual-work test additionally checks the motion/load
pairing under deformation.

\subsection{Supported connection scope}

First support tested rigid connections and appropriate free equilibrium
locations. The external-resultant equation map must account for constrained
degrees of freedom and reactions. Hinge-side rotational equations require
additional handling; silently applying a moment to the first matching equation
is unacceptable. Reject unsupported mappings until their assembly is verified.
Rebuild maps after topology or equation numbering changes.

\section{Coupled step lifecycle and rollback}
\label{sec:lifecycle}

\subsection{One physical step can contain many numerical trials}

Let $S_n$ and $A_n$ denote accepted structural and aerodynamic states at time
$t_n$. Outer coupling iterations search for a consistent pair at
$t_{n+1}=t_n+\Delta t$. They are not additional physical time steps.

\begin{enumerate}
\item Capture accepted structural/aerodynamic history and external-load state.
Prepare the structural step, basis motion, boundary conditions and equivalent
rates once for this attempted step.
\item Obtain the current structural trial geometry and velocities.
Restore accepted aerodynamic history, apply the trial geometry, and evaluate
global UVLM loads with wake advancement disabled.
\item Transfer dimensional forces/moments to the structural load buffer.
Hold these loads fixed in basis A during the inner structural Newton solve.
\item Evaluate UVLM again at the resulting structural iterate. Check changes
in state and load, and equilibrium with the updated loads. Inner Newton
convergence with lagged loads is not sufficient.
\item Repeat with the configured relaxation until outer criteria pass.
\item Accept the consistent pair, advance the wake and rotor phase once,
commit active wake counts, and save accepted outputs.
\item On failure, restore both disciplines, phase, clocks, load buffers and
history. Report failure; enable smaller-step retries only after validation.
\end{enumerate}

The existing \code{advance_uvlm_trial!} helper supports disabling wake
advancement, but its current default is to advance. The new coupling-facing
API must make a noncommitting evaluation the safe, explicit operation.

\subsection{Where AeroBeams needs generic additions}

Add a proposed \code{ExternalNodalResultants} container with dimensional
forces and moments in basis A. Apply its residual contribution after normal
element and special-node assembly in \code{assemble_system_arrays!}.
The buffer can initially follow the earlier blueprint's placement on
\code{Model}, with documented single-run mutable ownership and inclusion in
the bridge snapshot.

Expose generic prepare, solve, accept, and restore step operations around the
existing structural machinery. Factor code from the current lifecycle instead
of copying the whole time integrator into the bridge. Initial-state/rate
initialization must receive consistent external loads as well.

The current structural snapshot is not a complete coupled transaction. Audit
time/index values, reference-frame orientation, current and equivalent rates,
constraint and boundary-condition state, complementary aerodynamic variables,
solver status, load buffers, and saved histories. A failed step must not leave
a saved time entry or an advanced rotor phase.

\subsection{Partitioned first; monolithic later}

Frozen loads have no additional UVLM derivative inside the initial inner
structural Newton solve. Overall consistency is established by the outer
iteration. Strong added-mass or other difficult regimes may require better
relaxation or a stronger coupling method; no universal convergence guarantee
is assumed.

A later monolithic method needs global aerodynamic derivatives, including
cross-element/surface blocks. Current element-local aerodynamic derivatives
do not provide them. Validate finite-difference derivatives at fixed accepted
wake history before adopting analytical or AD-based derivatives. Scalar type
restrictions and the existing force-derivative limitations must be addressed
explicitly.

\section{Rotor modeling and analysis capability differences}

Prescribed rotor speed is the first target. Each rotor has its own speed,
shaft direction, and phase. The requested speed scaling remains a prescribed
operating policy; it is not a dynamic motor/shaft equation.

Native inertial components can represent rotor mass and nonspinning inertia,
but relative shaft spin produces additional angular momentum. Schematically,
\begin{equation}
 h_{\mathrm{spin}}=J_s\Omega e_s.
\end{equation}
Its time derivative contributes to the rotor/support angular-momentum balance.
The support-reaction sign must be derived from the chosen balance convention
and tested against the retained Chang small-angle gyroscopic matrix. Avoid
counting the same rigid-body inertia both in a point inertia and a custom
rotor component.

The validation sequence is: structural-only beam behavior, zero-spin rotor,
small-angle gyroscopic behavior, prescribed-spin coupled response, then
mesh/time/coupling convergence. A physical AeroBeams pylon and the Chang
two-coordinate modal pylon are different models. Retain both as explicitly
named alternatives until their intended comparison is established.

Variable-speed rotor dynamics, general hinges, physical mount damping, steady
trim, and eigen/Floquet analyses are later features. In particular, an
autonomous eigenproblem about a steady state and stability about a periodic
rotor state are not interchangeable.

\section{Moving-block damping, studies, and recorded results}

\subsection{Preserve the method selected by the user}

The current moving-block implementation is retained during migration:
selected positive peaks bound the fitting segment; the reference-sized block
occupies 25--50\% of that segment; blocks shift one sample; each block has its
mean removed; a rectangular FFT identifies the dominant non-DC component;
and a line is fitted to log amplitude versus block-start time:
\begin{equation}
 \log A_k\simeq a+\lambda t_k.
\end{equation}

Move the method into the studies layer through an explicit proposed
\code{MovingBlockDampingMethod} setting. Retire the temporary entry-local
wrapper only after the central runner reproduces its behavior. Preserve
pitch/yaw diagnostics, optional fixed-duration mode, band acceptance,
minimum peaks, fit quality, and estimator identity.

This is a location/API change, not permission to replace the estimator.
Offset sensitivity retained from positive-peak screening remains documented.
A detrended or multimode variant would require a distinct method/version.

\subsection{Frequency resolution and sampling}

For $N_b$ samples at spacing $\Delta t$, FFT bin spacing is
\begin{equation}
 \Delta f=\frac{1}{N_b\Delta t}.
\end{equation}
A convergence tolerance of 0.1~Hz does not establish that frequency estimates
are accurate to 0.1~Hz when the block resolution is coarser. Record
\code{frequencyResolutionHz} and retain the off-bin regression cases.
Bin spacing is not itself a statistical uncertainty bound.

Moving-block FFT requires uniform sampling. Keep dense lightweight responses
even if geometry snapshots are saved sparsely. An adaptive accepted time grid
must not be passed directly into this FFT. Initially restrict this damping
workflow to fixed-step responses; any later resampling must record and validate
its interpolation and effects on the estimate.

\subsection{Completion and scientific acceptance}

Result records must separately answer:
\begin{itemize}
\item Did the calculation complete the expected discrete time grid?
\item Did the structural and outer coupling solvers converge?
\item Is the damping estimate valid for the configured method and checks?
\item Did each convergence family pass, and did a combined refinement pass?
\end{itemize}

Structural mesh convergence becomes explicit after decoupling the meshes.
Wing aerodynamic convergence needs wing force/moment metrics; CT/CQ alone
cannot establish it, especially with disabled wing--propeller interaction.
Record requested and accepted grids, termination reason, residual histories,
units/bases, solver/method identities, and the resolved operating point.

\section{Compatibility, environments, and provenance}

\subsection{Old inputs and new names}

An input adapter translates legacy names once at construction. It validates
units and conflicts, then supplies one canonical configuration to the solver.
For example, a legacy degree-valued angle becomes a radian-valued canonical
angle; this is more than copying a number under a new key.

Reject conflicting old/new aliases instead of silently choosing one. Emit
deprecation notices once at the public boundary, not in time steps or panel
loops. Preserve current example paths and useful \code{run_chang} behavior
through thin wrappers. Remove wrappers only in a documented later release.

\subsection{Reproducible Julia environments}

The UVLM package, studies, plotting workflow, and bridge need declared
dependencies. The earlier review found that studies/tests used FFTW and Plots
through the global environment. Successful execution in that environment does
not establish a clean installation.

Use a dedicated study project, explicit local package development paths where
appropriate, and tests with \code{LOAD_PATH=["@","@stdlib"]}. Constructors
and imported modules must not activate a project or start a study as a side
effect. Executable scripts should guard their entry points.

\subsection{Existing aerodynamic selections cannot be silently resealed}

The current numerical fingerprint includes source paths and contents across
the backend, Chang implementation, and convergence machinery, plus environment
information. Moving or renaming source files may invalidate it even when the
formulas are preserved.

Keep archived selections and seals unchanged. They remain readable for
historical inspection/postprocessing with their original schema. To create a
new verified selection under migrated code, rerun or explicitly verify the
required cases under that source identity. Passing rename-regression tests
does not authorize replacing the old seal with a new one.

The new provenance design should distinguish aerodynamic, structural,
coupling, estimator, and output-schema identities with conservative dependency
manifests. Include indirect dependencies. Source identity must correspond to
the code actually loaded: fresh processes are the initial policy for verified
runs after edits. Cache keys must include every relevant input, including
speed/RPM, excitation, initial state, time grid, and tolerances.

\section{Implementation phases and acceptance gates}
\label{sec:phases}

The phase identifiers match the existing Markdown roadmap. Each phase ends
with migrated call sites, tests, documentation, and an explicit support status.
The first six phases provide the terminology/organization migration; actual
structural coupling follows.

\subsection{P0: capture the baseline}

Inventory public exports, entry points, configuration keys, output schemas,
and dependency assumptions. Save resolved current presets and small numerical
fixtures for rigid wing, propeller, interactions, and the Chang response.
Record versions and source identities. Retain the already recorded moving-block
tests as evidence for that method.

\textbf{Acceptance:} current supported workflows can be reproduced in fresh
processes; inputs and outputs are recorded; limitations are explicit.
\textbf{Dependency:} none. No renaming should precede a reproducible baseline.

\subsection{P1: isolate environments and unify entry behavior}

Declare study/test dependencies, keep plotting optional, guard executable
entry points, and centralize the public stage runner. Preserve the recently
implemented aeroelastic speed override and proportional RPM delegation.
Including a reusable file should not launch a full convergence study.

\textbf{Acceptance:} isolated environments load successfully; inclusion causes
no solve/output; all documented aeroelastic run paths resolve the same point.
\textbf{Dependency:} P0.

\subsection{P2: add a naming and constructor facade}

Implement the canonical types and constructor vocabulary over the existing
backend. Add legacy-key conversion, validation and deprecation behavior.
Separate named case builders from generic constructors. Document input,
derived and mutable fields, units and ownership.

\textbf{Acceptance:} old/new settings resolve identically, conflicts are
rejected, and baseline numerical fixtures match.
\textbf{Dependency:} P1.

\subsection{P3: separate model, state, workspace, and problem}

Classify the existing \code{System} storage and introduce the public
model/state/workspace facade. Implement explicit initialization, trial,
commit and rollback. Preserve one authoritative storage owner and isolate
independent runs.

\textbf{Acceptance:} repeated trials are deterministic; A--B--A trials recover
the same A result; rejected steps leave no history change; duplicate commit
is rejected; saved results do not alias live workspaces.
\textbf{Dependency:} P2.

\subsection{P4: reorganize results, damping, and studies}

Introduce structured status/results and history policies. Move generic study
machinery out of the Chang example. Select moving-block damping through an
explicit estimator configuration, replacing entry-specific dispatch only
after equivalence is proven.

\textbf{Acceptance:} current estimator and diagnostic fixtures match;
incomplete or unconverged runs cannot be accepted as converged studies;
lightweight responses are available without full geometry history.
\textbf{Dependency:} P3.

\subsection{P5: version outputs and migrate documentation}

Implement conservative stage-specific identities and schema-versioned readers.
Update the API/library guide, editable examples, migration dictionary, and
cache policy. Keep original aerodynamic selections intact.

\textbf{Acceptance:} legacy results remain readable; stale caches are rejected;
new verified outputs refer to the correct source/environment and input point.
\textbf{Dependency:} P4. This completes the first recommended approval scope.

\subsection{P6: add generic AeroBeams step and load interfaces}

Implement external dimensional resultants and structural prepare/solve/accept/
restore operations. Audit initial-state initialization and snapshot coverage.
Ensure the default behavior without external loads is preserved.

\textbf{Acceptance:} existing structural tests pass; known loads have correct
signs, scaling and reactions; forced rejection restores all mutated state.
\textbf{Dependency:} P0 and the transaction contract established in P3.

\subsection{P7: implement the motion/load bridge}

Create the separate bridge package and mapping objects. Implement explicit
frames, independent meshes, geometry/velocity reconstruction, dimensional load
transfer, attachment stencils, and supported connection validation.

\textbf{Acceptance:} rigid-motion, resultant force/moment, virtual-work,
nonmatching-mesh and moving-frame tests pass.
\textbf{Dependencies:} P5 and P6.

\subsection{P8: couple fixed-step dynamics}

Implement the wrapper problem and outer iteration using frozen inner-solve
loads. Check the final consistent state/load pair. Support initialization,
structured failure and exactly one accepted wake commit.

\textbf{Acceptance:} rigid prescribed motion, zero-aerodynamic-load limit,
repeated trials and a small flexible-wing coupled benchmark pass.
\textbf{Dependency:} P7.

\subsection{P9: validate the Chang rotor case with AeroBeams}

Build the AeroBeams structural counterpart separately from the retained Chang
model. Validate mass, static response, modes, attachments, and rotor spin
reaction before full aeroelastic comparisons. Start with zero physical
structural damping unless a supported damping component is available.

\textbf{Acceptance:} structural-only correspondence, zero-spin and small-angle
gyroscopic checks, then coupled mesh/time/iteration convergence.
\textbf{Dependency:} P8.

\subsection{P10: optimize based on measurements}

Apply history controls, buffer reuse, explicit matrix cache invalidation and
independent interaction-group solves where supported. Profile compilation,
warm numerical work, allocation and I/O separately.

\textbf{Acceptance:} fixture loads/states remain equivalent; measured memory
and warm-run performance justify each change.
\textbf{Dependencies:} the relevant ownership/coupling phases P3--P9.

\subsection{P11: extend analysis capabilities}

Add adaptive coupling and validated response resampling, steady/trim variants,
global aerodynamic derivatives, eigen/periodic stability support, broader
connections, and rotor models only as separate numerical tasks.

\textbf{Acceptance:} each feature has its own benchmark and limitation record;
unsupported modes are rejected until implemented.
\textbf{Dependencies:} P8/P9 and the specific feature design.

\subsection{Ordering and review units}

The primary path is
\begin{equation*}
 \mathrm{P0}\to\mathrm{P1}\to\mathrm{P2}\to\mathrm{P3}\to
 \mathrm{P4}\to\mathrm{P5}\to\mathrm{P7}\to\mathrm{P8}\to\mathrm{P9},
\end{equation*}
with P6 supplying the structural interface required by P7. Performance work
begins when the affected ownership contracts are stable. A calendar estimate
should follow P0 measurements rather than being inferred from the number of
files to rename.

Each reviewable implementation change should state: the concrete behavior
before/after, files migrated, numerical invariants preserved, validation run,
output/provenance impact, and remaining unsupported cases.

\section{File-by-file migration responsibilities}

Paths below are relative to \code{lib/WingPropellerUVLM}, except rows marked
root. Grouped filenames describe migration targets, not new files already
created.

\begingroup\small
\begin{longtable}{@{}L{0.38\linewidth}L{0.54\linewidth}@{}}
\toprule
Current area & Planned action\\
\midrule\endhead
\code{src/WingPropellerUVLM.jl} & Add public includes/exports incrementally;
retain package identity and compatibility boundaries.\\
\code{src/backend/system.jl}, \code{src/UVLMState.jl} & Inventory fields;
split ownership through a facade; formalize transactions and alias rules.\\
Backend panel, geometry, circulation, induced, analyses and wake files &
Adapt first, preserve formulas/chronology, then rename internals if useful.\\
Backend nearfield, nearfield postprocessing, reference and freestream files &
Document normalization and physical application points; expose in-place
dimensional load buffers.\\
Legacy nearfield, nonlinear, rotor import, farfield, stability and visualization &
Inventory exports and supported behavior. Retain unported workflows behind
compatibility interfaces; do not claim new validation.\\
\code{src/wing_propeller} & Separate generic grid/motion operations from
tabulated case data; retain ordering, interpolation, twist and spin signs.\\
\code{examples/.../src/ChangAeroelastic.jl}, configuration and parameters &
Thin benchmark builder plus input translation; generic construction moves
to the package.\\
Chang structural matrices/model & Retain the linear benchmark and add a
separate native AeroBeams builder in P9.\\
Chang workspaces, UVLM coupling and simulation & Extract generic storage/loop
while preserving the Chang-specific kinematic adapter and virtual-work tests.\\
\code{src/aeroelastic} & Keep generalized-alpha for the benchmark; separate
excitation from integration and avoid advancing AeroBeams with it.\\
Chang postprocessing, plots and animation & Structured outcomes, explicit
writers, optional expensive geometry tracking.\\
Convergence entries, adapters, selection/options and metrics & Central study
runner, explicit estimator, legacy input/output readers, conservative provenance.\\
Root \code{src/Model.jl}, \code{src/Problem.jl},
\code{src/SystemSolver.jl}, proposed external resultants & Add generic load
and step hooks; preserve existing uncoupled solver behavior.\\
Root \code{src/Core.jl} & Audit residual, rotation and scaling contracts;
change only necessary generic interfaces.\\
Project files, tests and docs & Declare dependencies; isolate environments;
update public API and executable examples with each phase.\\
\bottomrule
\end{longtable}
\endgroup

\section{Validation plan and completion criteria}

\begingroup\small
\begin{longtable}{@{}L{0.24\linewidth}L{0.68\linewidth}@{}}
\toprule
Group & Required checks\\
\midrule\endhead
Input/API & Old/new aliases; units; unknown/conflicting keys; defaults;
65-to-85 m/s override with RPM scaling; include without execution.\\
Environment & UVLM alone, studies, optional plots and bridge load with declared
dependencies and without the global environment.\\
Aerodynamic regression & Existing core/panel/force tests; rigid wing, rotor,
interacting case; circulation, rates, dimensional forces and wake geometry.\\
State transactions & Repeated trials, A--B--A sequence, failed step, duplicate
commit, wake capacity, phase rollback and result-array isolation.\\
Structural integration & Existing AeroBeams suite, known applied force/moment,
dynamic loads, reactions and supported special-node mappings.\\
Transfer & Independent meshes, nonzero rotations, offsets, moving frame,
resultant forces/moments and virtual work.\\
Coupled limits & Zero aerodynamic loads; prescribed rigid structure; wing
without propeller; disabled cross-group interaction; zero-spin rotor.\\
Damping & Reference moving-block comparison; growth/decay; off-bin and
out-of-band frequencies; offsets/noise/mode switching; sampling and diagnostics.\\
Study acceptance & Actual grid completion, inner/outer convergence, finite
state limits, valid metrics, individual and combined refinements.\\
Provenance & Old seals unchanged; overridden reference recorded correctly;
changed source, method or inputs prevent invalid cache reuse.\\
Performance & Warm timings and allocated bytes per major operation; compilation
and I/O separated; storage growth and run isolation measured.\\
\bottomrule
\end{longtable}
\endgroup

The earlier review recorded 369 existing assertions and 13 additional
smoke/virtual-work assertions. The subsequent moving-block work recorded
124 targeted checks. These are historical results from different runs with
overlapping coverage, not a new summed test total or proof of a future bridge.
The short 38-step smoke does not validate a six-second damping response.

For pure naming changes, require identical resolved inputs and deterministic
outputs. Where evaluation order changes, define scaled tolerances from the
fixtures and show that discrepancies are well below study thresholds.
Do not relax tests simply to accept a new result. For a structural formulation
change, compare converged physical observables and modeling assumptions rather
than requiring bitwise equality with the old integrator.

\section{Efficiency improvements enabled by the structure}

The first efficiency improvement is predictable history storage. Dense
response samples are inexpensive compared with saving all surface panels and
wake geometry every step. Make heavy histories opt-in with a stride, while
preserving the response grid needed for damping.

Next reuse geometry, force, transfer, and snapshot buffers. Cache the retained
Chang generalized-alpha effective matrix only while its M/C/K, time step and
integration parameters are unchanged. Gyroscopic terms mean a symmetric
positive-definite factorization cannot be assumed. AeroBeams retains its own
Jacobian update policy.

Aerodynamic matrix reuse must respond to geometry, core, interaction and
relevant flow/time changes. Panel-count equality alone is not a valid cache
key. Disconnected interaction groups can be solved independently only while
preserving interactions inside each group, including blades of one rotor.

Measure geometry, AIC assembly/factorization, force calculation, wake
convection, transfer, structural solve, snapshots, estimator and output
separately. Process-based case concurrency needs measured memory limits and
independent output directories. Wake coarsening, fast multipole methods and
new integration formulas are numerical-method changes outside the initial
terminology migration.

\section{Risks, decisions, and approval checkpoint}
\label{sec:approval}

\begin{longtable}{@{}L{0.30\linewidth}L{0.62\linewidth}@{}}
\toprule
Main risk & Planned control\\
\midrule\endhead
A new name hides changed units or physics & Explicit semantic dictionary,
constructor conversion tests, and separate formulation phases.\\
Existing aerodynamic results stop verifying after file moves & Preserve
archives/seals, version identities, and reverify under migrated source.\\
Wake or rotor phase advances during repeated trials & Explicit transaction
operations, single commit token and forced-rejection tests.\\
Wrong force or moment reaches the beam equations & Dimensional frame-aware
buffers, scaling/reaction tests and virtual-work verification.\\
The structural replacement changes the benchmark & Separate structural-only
calibration and rotor gyroscopic validation before coupled comparisons.\\
Public APIs become inconsistent during migration & One canonical constructor/
runner path, boundary adapters and scheduled compatibility retirement.\\
Scopes grow before a usable result exists & Deliver P0--P5 first and use
fixed-step, prescribed-spin coupling as the first bridge target.\\
\bottomrule
\end{longtable}

\textbf{Recommended approval statement for the first implementation scope:}

\begin{quote}
Proceed with P0--P5 to introduce AeroBeams-style terminology, constructors,
problem/state/workspace organization, study/result handling and provenance,
while retaining the current UVLM numerical backend and Chang benchmark.
Preserve the moving-block method, speed/RPM policy and legacy readers.
\end{quote}

The user may instead approve the entire staged roadmap or request revisions
to names, package boundaries, compatibility duration, or priorities. Approval
of the document's direction should state whether it covers P0--P5 only or
also the P6--P9 structural coupling work. Advanced capabilities remain gated
by their own validation requirements.

No migration implementation is performed as part of this document delivery.

\appendix
\section{Source basis and relationship to the earlier documents}

The local source is authoritative for the proposed interface. Relevant files
in the root package are \code{src/Model.jl}, \code{src/Beam.jl},
\code{src/Element.jl}, \code{src/Problem.jl}, \code{src/SystemSolver.jl},
\code{src/Core.jl}, \code{src/Utilities.jl}, \code{src/UnitsSystem.jl},
\code{src/AeroSolver.jl}, and \code{src/AeroSurface.jl}.

Relevant UVLM sources include \code{src/backend/system.jl},
\code{src/UVLMState.jl}, \code{src/aeroelastic/GeneralizedAlpha.jl},
the Chang model/simulation/coupling files, and the two convergence entry
scripts and moving-block module under their current example directory.

The detailed Markdown worklist is
\code{docs/plans/UVLM_AEROBEAMS_MIGRATION.md}. Historical correctness evidence
is in \code{docs/reviews/uvlm_aerobeams_2026-09-14/REVIEW.md} and its associated
smoke/diagnostic files. The existing mathematical bridge blueprint is
\code{docs/src/wing-propeller-uvlm-coupling.md}; its prototype interfaces require
review before implementation. This LaTeX proposal explains and retains the
P0--P11 sequence while expanding the meaning of the principal differences.

The earlier review predates the current moving-block entry and corrected
nested speed-override delegation. Those changes are treated as existing
baseline behavior in this proposal; unrelated earlier findings are not
assumed resolved.

For general project context, see the
\href{https://github.com/luizpancini/AeroBeams.jl}{official AeroBeams repository}.
For qualified method extension and module naming, see the
\href{https://docs.julialang.org/en/v1/manual/modules/}{Julia modules manual}.
These references supplement the inspected checkout; proposed UVLM constructors
must not be mistaken for APIs documented upstream.

\section{Compiling and reviewing this document}

This source is self-contained and uses standard LaTeX packages. No Julia
code in the listings is executed during compilation. No external images,
bibliography database, or shell escape are required.

With a local LaTeX distribution, run from the source directory:
\begin{lstlisting}
pdflatex -interaction=nonstopmode -halt-on-error UVLM_AEROBEAMS_MIGRATION_PROPOSAL.tex
pdflatex -interaction=nonstopmode -halt-on-error UVLM_AEROBEAMS_MIGRATION_PROPOSAL.tex
\end{lstlisting}
The second pass resolves the table of contents and internal references.
Alternatively use \code{latexmk -pdf}, or upload the single source to Overleaf
and select pdfLaTeX. The current authoring environment did not provide a
LaTeX compiler, so PDF compilation and rendered-page inspection were not
performed here. Static source checks are not a substitute for that review.

\end{document}
